\documentclass[12pt,letterpaper]{article}
\usepackage[margin=1in]{geometry}
\usepackage{amsmath,amssymb,amsthm,blkarray,braket,caption,enumitem,fancyvrb,graphicx,hyperref,relsize,verbatim}
\usepackage[numbers]{natbib}
\usepackage[section]{placeins}
\graphicspath{{./images/}}
\bibliographystyle{apalike}

\begin{document}

\title{Chapter 1.7 - Group Actions}
\author{Emil Marinov}
\maketitle

\section{Problem 18}

\textit{Let $H$ be a group acting on set $A$.  Prove that the relation $\sim$ on $A$ defined by}

\begin{equation}
   \begin{aligned}
      a \sim b \quad \text{iff } a=hb \quad \text{for some } h \in H
   \end{aligned}
\end{equation}

\textit{is an equivalence relation.
(For each $x \in A$ the equivalence class of $x$ under $\sim$ is called the \textbf{orbit} of $x$ under 
the action of $H$.
The orbits under the action of $H$ partition the set $A$.)} \\
\\
To show an equivalence, reflexivity, symmetry, and transitivity must be shown. \\
\textbf{Reflexivity} states that $(a, a) \in \rho \forall a$, where $\rho$ is the relation.

\begin{equation}
   \begin{aligned}
      a = ea
   \end{aligned}
\end{equation}

As $e \in H$, there exists an $h = e \in H$ such that $a = ha$.  Thus, $a \sim a$.

\textbf{Symmetry.} The relation on $a$ and $b$ is symmetric iff $(a,b) \in \rho \implies (b,a) \in \rho$.
Let $a = hb$.

\begin{equation}
   \begin{aligned}
      a = hb \\
      h^{-1} a = (h^{-1} h)b \\
      h^{-1} a = b
   \end{aligned}
\end{equation}

$h^{-1} \in H$ and there exists a $g = h^{-1} \in H$ such that $b = ga$.
Thus symmetry has been shown.

\textbf{Transitivity.} The relation $\sim$ is transitivity iff $(x,y) \in \rho, (y,z) \in \rho \implies (x,z) \in \rho$.
Let $a = h_1b$ and $b = h_2c$, for $h_1, h_2 \in H$.

\begin{equation}
   \begin{aligned}
      a = h_1b \\
      a = h_1(h_2c) \\
      a = (h_1 h_2)c \\
      a = h_3c && \text{Let } h_3 = h_1 h_2 \in H \\
      a \sim c
   \end{aligned}
\end{equation}

The group operation $h_1 h_2$ yields an element $h_3 \in H$ due to closure.

The relation $~$ on set $A$ defined by $a \sim b \quad \text{iff } a=hb \quad \text{for some } h \in H$
has been shown to be reflexive, symmetric, and transitive.
Thus, it has been proven to be an equivalence relation. $\square$ \\
\\
The same goes for $y \in A$, or any other element in A.
The orbits of different elements $x \in A$ under the action of $H$ are disjoint, as indicated by the equivalence relation.
If they were not disjoint, then the equivalence relation indicates they share ALL elements together 
i.e. being the same orbit.


\section{Problem 20}
\textit{Let $G$ be a group and let $G$ act upon itself by left conjugation, so each $g \in G$ maps $G$ to $G$ by \\
\begin{equation}
   \begin{aligned}
      x \to gxg^{-1}
   \end{aligned}
\end{equation}
For fixed $g \in G$, prove that conjugation by $g$ is an isomorphism from $G$ onto itself (i.e., is an \textbf{automorphism} of $G$
- cf. Exercise 20, Section 6).
Deduce that $x$ and $gxg^{-1}$ have the same order for all $x$ in $G$ and that for any subset $A$ of $G$, 
$|A| = |gAg^{-1}|$ (here $gAg^{-1} | a \in A$).} \\
\\
To show that the left conjugation on $G$ is an isomorphism, it must be shown that 1) it is a homomorphism, 2) it is injective, 
and 3) it is surjective.

\textbf{Show Homomorphic} \\

\begin{equation}
   \begin{aligned}
      \varphi(xy) = g(xy)g^{-1} \\
         = gxyg^{-1} && \text{By associativity.} \\
         = gxg^{-1}gyg^{-1} \\
         = (gxg^{-1})(gyg^{-1}) \\
         = \varphi(x) \varphi(y)
   \end{aligned}
\end{equation}

\textbf{Show Injective} \\

Here, it is needed to be shown that $\varphi(x) = \varphi(y) \implies x = y$.

\begin{equation}
   \begin{aligned}
      \varphi(x) = gxg^{-1} = gyg^{-1} = \varphi(y) \\
      gxg^{-1}g = gyg^{-1}g \\
      gx = gy \\
      g^{-1}gx = g^{-1}gy \\
      x = y
   \end{aligned}
\end{equation}

\textbf{Show Surjective}

To show surjectivity, it must be shown that each element from the codomain has an element in the domain 
that maps to it: $\forall y \in G, \exists x \in G \quad \text{s.t.} \varphi(x) = y$.

\begin{equation}
   \begin{aligned}
      y = \varphi(x) = gxg^{-1} \\
      \implies x = g^{-1} y g
   \end{aligned}
\end{equation}

This holds for all $g \in G$.  Thus, it is surjective.

\textbf{Order} \\

Let $n$ be the order of $x$, i.e. $x^n = 1$.
Let $m$ be the order of $\varphi(x)$, i.e. $\varphi(x)^m = 1$.

\begin{equation}
   \begin{aligned}
      \varphi(x)^m = (gxg^{-1})^m 
      = gx^mg^{-1} = e \\
      x^m = g^{-1}eg = g^{-1}g = e
   \end{aligned}
\end{equation}

If $x^n = e$, then $gx^ng^{-1} = e$.  This implies $|\varphi(x)| = m \le n$.
If $gx^mg^{-1} = e$, then $x^m = e$.  This implies $|x| = n \le m$.

Because $m \le n$ and $n \le m$, $m=n$.
Thus, the order of the automorphism is the same as that for the input group element.

\textbf{Order 2}

$gAg^{-1}$ is bijective with $A$.  Thus, $|A| = |gAg^{-1}|$.

\end{document}